\documentclass[../main.tex]{subfiles}

\begin{document}

\chapter{Introduction}
\label{chap:intro}

% ============================================================
\section{Motivation and Context}
\label{sec:motivation}
% ============================================================

The quality of training data is widely recognised as one of the most critical determinants of the performance and reliability of machine learning models.  Real-world datasets are seldom pristine: they are pervasively affected by a wide range of imperfections (including some of which we will deeply explore in the themes of this thesis like: missing values, mislabelled instances, duplicate records, noisy measurements, and statistical outliers).  These defects may arise from heterogeneous sources (sensor malfunction, human annotation error, data integration across disparate systems, etc..).  When such corrupted data are consumed by learning algorithms, the resulting models may exhibit degraded generalisation, biased predictions, or, in safety-critical domains such as cybersecurity and healthcare, potentially harmful decision-making.

Within the paradigm of Data-Centric AI, the focus of the deep learning pipeline widely shifts from model architecture design to systematic improvement and curation of the training data.  Central to this paradigm is the ability to quantify the impact of specific data quality issues on downstream model performance.  This capability, however, presupposes the availability of tools that can introduce controlled, reproducible perturbations into a dataset while leaving all other aspects of the data intact.
Such tools enable researchers to answer questions of the form: "How much does model predicting perfomances degrade when $x$\% of a given feature's values are replaced by noise?"

\texttt{PuckTrick}\supercite{pucktrick_docs} is an open-source Python library developed at the University of Milano-Bicocca that addresses precisely these needs. \texttt{PuckTrick} provides a unified, programmatic interface for the controlled introduction of five fundamental categories of data corruption into tabular datasets intended for machine learning applications: \emph{duplicated records}, \emph{missing values}, \emph{noisy values}, \emph{outliers}, and \emph{label inconsistencies}.  Each corruption type is parametrized witch the choice of a target column and the desired percentage of affected rows, enabling fine-grained, reproducible experimental designs in which individual data quality dimensions can be isolated and studied independently.

In its original form, PuckTrick operates exclusively on Pandas DataFrames (a widely used in-memory data structure for tabular data in the Python ecosystem).  While Pandas is an excellent choice for small-to-medium datasets that fit comfortably in a single machine's memory, it becomes a bottleneck when the dataset scales to millions or tens of millions of rows.  Modern cybersecurity, IoT, and financial datasets routinely exceed the single-node memory capacity, and the scientific community increasingly relies on distributed computing frameworks to process them.

% ============================================================
\newpage
\section{The Cybersecurity Application Domain}
\label{sec:cybersecurity-domain}
% ============================================================

The empirical component of this thesis is situated within the domain of \emph{network intrusion detection}: a field that offers a particularly compelling setting for studying the interplay between data quality and deep learning performance.
By it's nature, this setting offers wide range of data that can be used in a deep learning context for developing an automatic intrusion detection system (based on deep learning): a single day of traffic capture at an institutional network can produce millions of flow records, each described by dozens of statistical features (packet counts, byte volumes, inter-arrival times, TCP flag distributions, etc.). Processing and corrupting datasets of this magnitude requires distributed computing capabilities (precisely the motivation for extending PuckTrick to Apache Spark).
A \emph{Network Intrusion Detection System} (NIDS) is tasked with automatically classifying each observed network flow as either benign or malicious and, in the latter case, identifying the specific type of attack being carried out.(such as: brute-force credential guessing, Denial-of-Service (DoS) and Distributed Denial-of-Service (DDoS) floods, web-application exploits (SQL injection, cross-site scripting), botnet command-and-control communications, and stealthy infiltration campaigns).

The specific dataset selected for this study is the \textbf{CSE-CIC-IDS2018}\supercite{ids2018} (\emph{Canadian Institute for Cybersecurity Intrusion Detection System 2018}), a publicly available benchmark produced by the University of New Brunswick and distributed on the Kaggle platform.
The dataset was collected over a ten-day capture window in 2018 inside a controlled cyber-range infrastructure: 50 attack machines launched a curated battery of attacks against a victim network comprising 420 machines and 30 servers, while a profile generator replayed realistic benign traffic drawn from a usage model encompassing HTTP, HTTPS, FTP, SSH, and email activity.
All bi-directional network flows were extracted by the \textit{CICFlowMeter} tool, which computes 80 statistical features.

In its unified form the dataset comprises over \textbf{16 million network-flow instances}. While lacking presence of any zero-day attacks, as all the network flows were simulated, having this large amount of network flows compels making it one of the largest and most realistic publicly available benchmarks for intrusion detection research.  A detailed statistical characterization of the dataset  is provided in Chapter~\ref{chap:dataset}.

The choice of this domain is deliberate: by studying the effects of controlled noise injection on deep learning models that must discriminate between benign and malicious network traffic, this thesis bridges two active research areas: \emph{data quality for machine learning} and \emph{AI-driven cybersecurity}. This will provide actionable insights into how resilient modern deep learning architectures are when the data on which they are trained exhibit quality degradation of known type and magnitude. 

% ============================================================
\newpage
\section{Objectives of the Thesis}
\label{sec:objectives}
% ============================================================

This thesis pursues two interrelated objectives:

\paragraph{Objective 1: Extending PuckTrick to Apache Spark.}
The primary engineering contribution of this work is the extension of the PuckTrick library to support \textbf{Apache Spark} (via PySpark) as an alternative processing backend alongside Pandas.  The work was conducted within the \texttt{Spark\_Conversion} branch of the PuckTrick repository in collaboration with Daniel Satriano, under the supervision of Professor Andrea Maurino.



\paragraph{Objective 2: Evaluating the impact of controlled noise on deep learning models for cybersecurity.}
The second objective is an empirical investigation into the effects of PuckTrick-generated data corruption on the performance of deep learning models trained for \emph{network intrusion detection}, the cybersecurity application domain introduced in Section~\ref{sec:cybersecurity-domain}.  

The investigation operates at two levels of granularity:
\begin{itemize}
    \item At the \emph{binary classification} level: the models are asked to discriminate between benign and malicious flows: a task analogous to anomaly detection, where any deviation from normal traffic triggers an alert.
    \item   At the\emph{multi-class classification} level, the models must additionally identify the specific type of attack, a substantially harder task that requires the classifier to learn fine-grained distinctions between attack profiles that often share similar flow-level statistics. 
\end{itemize}

Studying both levels simultaneously is essential because data corruption may have asymmetric effects: noise that does not materially degrade binary detection (benign vs. malicious) may nonetheless collapse the classifier's ability to distinguish among attack subtypes, or vice versa. In order to do so two state-of-the-art deep learning architectures were selected for the evaluation:

\begin{itemize}
    \item \textbf{TabNet}: a sequential attention-based architecture specifically designed for tabular data.  TabNet performs instance-wise feature selection through learnable sparse attention masks at each decision step, producing natively interpretable feature importance scores.  It was configured as a \emph{multi-task classifier}, simultaneously predicting both the binary (benign/malicious) and the fine-grained multi-class (specific attack type) labels.
    \item \textbf{CNN-LSTM model}: an ad-hoc model created on the basis of another study that aims to use our same dataset and try to improve a regular CNN architecture, using LSTMs blocks inside the model. This is done to help the model learn based on data given to it in time windows (rendering CNN capable of learning on time series such as this dataset), in spite of regular CNN tabular approach
\end{itemize}

As later specified in Chapter \ref{feature-relevance} and Chapter \ref{chap:5-Empiric-analysis-models} two main branches of experiments were developed:
\begin{itemize}
    \item \red{Experiment A:} \textbf{Experiments conducted starting from a reduced dataset}: these experiments were performed on a dataset preprocessed to remove both redundant features (features exhibiting high correlation in between each other) and zero-variance features. The primary objective was to identify noise injection strategies capable of consistently improving model performance. This phase serves as a proof of concept, demonstrating that the injection of controlled noise into the original dataset can enhance the effectiveness of deep learning-based Intrusion Detection Systems (IDSs).

    \item \red{Experiment B:} \textbf{Experiments conducted starting from full dataset}: in contrast to the previous setting, a broader set of experiments was carried out using the complete, unfiltered dataset. This experimental phase is intended to address the more theoretical aspects of the thesis. Its main goal is to provide an in-depth analysis of how the selected models respond to noise introduced through the \texttt{Pucktrick} library.
\end{itemize}


% ============================================================
\newpage
\section{The library implementation protocol}
% ============================================================
Before tackling any emppyrical experiment, this thesis will discuss the implementation of the \texttt{Pucktrick} library update. The central design principle guiding the extension was \emph{functional equivalence}: for every corruption method offered by \texttt{PuckTrick} in it's Spark implementation, it must've accepted the same parametrization and produce statistically equivalent output distributions in the newly implemented spark counterpart. Doing so the experiments conducted on a single-machine using Pandas backend would yield results that are directly comparable with those obtained in a distributed Spark cluster.  This requirement posed non-trivial engineering challenges, because Pandas and Spark embody fundamentally different computation models, and many of PuckTrick's corruption algorithms rely on random sampling, in-place mutation, and index-based row selection: operations that do not have direct analogues in the Spark API.

The implementation involved:
\begin{enumerate}
    \item A systematic audit of all existing PuckTrick corruption methods to
          catalogue their Pandas-specific dependencies.
    \item The design of a \emph{dual-backend architecture} in which each
          corruption module detects the type of the input DataFrame (Pandas
          \texttt{pd.DataFrame} or PySpark \texttt{pyspark.sql.DataFrame})
          and dispatches to the appropriate implementation.
    \item The re-implementation of each corruption method using PySpark's
          SQL-based DataFrame API, paying particular attention to
          pseudorandom-number generation with seed control, monotonically
          increasing identifiers, and window functions to simulate
          index-based operations.
    \item An extensive validation campaign comparing the outputs of the Pandas
          and Spark backends across all corruption types and all supported data
          types (continuous, discrete, binary, categorical string,
          categorical integer),
    \item Performance benchmarking on a multi-node Spark cluster to evaluate
          the distributed scalability of the extended library.
\end{enumerate}

% ============================================================
\newpage
\section{The Experimental Protocol}
% ============================================================

\red{QUESTO CAPITOLETTO DI EXPERIMENTAL PROTOCOL ANDREBBE RIVISTO TUTTO. IO SUGGERIREI ANCHE DI TOGLIERLO DALL'INTRODUZIONE E METTERLO PIUTTOSTO NELLA PARTE DI ESPERIMENTI}

The experimental protocol is structured as a multi-stage pipeline:
\begin{enumerate}
    \item  \textbf{Dataset analysis and preparation.} Firstly the entierty of the dataset  is put through an analysis and preparation pipeline. This will remove from the dataset any anomaly such as presence of null and inf values, as such as removal of 
        of low variance columns and reduntant ones (for more information on full pipeline visit Chapter~\ref{chap:dataset}). Also at the end of this pipeline there will be 
        an analys of fragility of remaininng features,  which classifies each feature as high,
        medium, or low fragility based on its skewness and kurtosis. This will allow 
        to effectively test noise injection into the dataset. 
    \item \textbf{Hyperparameter optimisation on a reduced dataset.}
        A stratified 0.1\% sample of the full dataset is used to conduct
        Bayesian hyperparameter optimisation via Optuna
        (Tree-structured Parzen Estimator sampler) for both TabNet and CNN-LSTM.
        The best hyperparameter configurations are persisted for subsequent
        reuse. Also, using best models from hyperparameters tuning, the attention-based feature importance scores of TabNet and the Variable Selection Network weights of TFT are extracted and intersected; this creates a composed list of approximately 10 most important features that are critical for both architectures.
    \item \textbf{Baseline and Feature importance analysis and selection.}
        The orginal dataset is then brought back scaled to 10\% of its original size (approximately
        1\,600\,000 instances). Using previously extracted hyperparameters and top 10 most
        important features 2 models (one for each model) are trained as baseline for subsequent
        experiments.
    \item \textbf{Feature fragility cross-referencing.}
          The feature importance rankings are then cross-referenced with the
          fragility analysis performed during dataset exploration
          (Chapter~\ref{chap:dataset}).  
          The final set of 5 target features for corruption comprises those that are
          simultaneously \emph{highly important} and \emph{highly fragile}: a
          combination that maximises the expected impact of noise injection on
          model performance.
    \item \textbf{Systematic noise injection and retraining.}
          Retaining same dataset as previous, Each of the five target features is then corrupted independently using four of PuckTrick's five corruption methods (missing, noise, outliers,
          labels), yielding $5 \times 4 = 20$ corrupted dataset variants.  The
          fifth method, \emph{duplicated}, is row-level and therefore applied
          once globally on an arbitrary feature.  Each variant is used to
          retrain both TabNet and TFT, producing a total of
          $(5 \times 4 + 1) \times 2 = 42$ model evaluations.
    \item \textbf{Statistical robustness via repeated trials.}
          The entire corruption-and-retraining loop is repeated 10 times with
          different random seeds to establish confidence intervals around the
          observed performance deltas, ensuring that the results are
          statistically robust and not an artefact of a particular random
          split or noise realisation.
\end{enumerate}

% ============================================================
\newpage
\section{Technologies and Frameworks}
\label{sec:technologies}
% ============================================================

This section provides a concise overview of the principal technologies and frameworks employed throughout this thesis, situating each within the broader landscape of data engineering and deep learning.

% ------------------------------------------------------------
\subsection{Apache Spark and PySpark}
\label{sec:tech-spark}
% ------------------------------------------------------------

Apache Spark is a unified analytics engine for large-scale data processing developed at the AMPLab of UC Berkeley and subsequently donated to the Apache Software Foundation.  Spark extends the MapReduce paradigm by introducing an in-memory computation model based on Resilient Distributed Datasets (immutable, fault-tolerant collections of records partitioned across a cluster's nodes).  Subsequent releases introduced the DataFrame API (a distributed analogue of the Pandas DataFrame with a SQL-like query optimizer, Catalyst) and the Dataset API (offering compile-time type safety on the JVM).

PySpark is the Python binding to Spark, enabling Python-native code to orchestrate distributed computation on a Spark cluster. In the context of this thesis, PySpark provides the execution substrate for the Spark-based corruption methods implemented in PuckTrick: each corruption operation is expressed as a sequence of Spark DataFrame transformations (column selections, filters, joins, window functions, and user-defined functions)  that are compiled by Catalyst into an optimized physical execution plan and distributed across the cluster's executors.

The Spark cluster used for benchmarking and experimental runs in this work
consists of a driver node and four executor nodes, each equipped with 4 cores
and 13\,GB of memory, for a total distributed capacity of 16 cores and 52\,GB
of RAM.


% ------------------------------------------------------------
\subsection{Pandas}
\label{sec:tech-pandas}
% ------------------------------------------------------------

Pandas is the standard library for, in-memory, tabular data 
manipulation in the Python data science ecosystem.  Its
core abstraction, the \texttt{DataFrame}, is a two-dimensional labelled data
structure with heterogeneous column types, supporting rich indexing, grouping,
merging, and reshaping operations.  Pandas operates eagerly and mutably on a
single machine: all data must reside in RAM, and transformations are executed
immediately upon invocation.

% ------------------------------------------------------------
\newpage
\subsection{TabNet}
\label{sec:tech-tabnet}
% ------------------------------------------------------------

TabNet~\supercite{arik2021tabnet} is a deep learning architecture explicitly designed for tabular data, proposed by Arik and Pfister at Google Cloud AI.  Its distinctive mechanism is \emph{sequential attention}: at each of $N_{\text{steps}}$ decision steps, a learnable attention mask selects a sparse subset of input features to process, and the outputs of successive steps are aggregated to produce the final prediction.

A more detailed explanation on how the model function is given in Chapter \ref{Chapter: model tabnet}


In this work, TabNet is instantiated as a \texttt{TabNetMultiTaskClassifier}
from the \texttt{pytorch-tabnet} library~\supercite{dreamquark2020tabnet}, with two
prediction heads sharing a common encoder: one for the binary label and one for the multi-class label.


% ------------------------------------------------------------
\subsection{Optuna}
\label{sec:tech-optuna}
% ------------------------------------------------------------

Optuna~\supercite{akiba2019optuna} is an automatic hyperparameter optimization
framework that implements efficient search strategies including the
\emph{Tree-structured Parzen Estimator} (TPE) and \emph{Hyperband} pruning.
It provides a \emph{define-by-run} API in which the search space is specified
programmatically inside the objective function, enabling dynamic space
definitions that depend on the values of previously sampled hyperparameters.

In this work, Optuna is used to tune both TabNet and CNN-LSTM.  The TPE sampler is
employed with early-stopping callbacks and configurable trial budgets (10
trials for local development, 20 for cluster-based runs).  Optuna's built-in
visualisation tools (parameter importance plots, optimization history) are used
to analyse the sensitivity of model performance to individual hyperparameters.

% ------------------------------------------------------------
\subsection{Scikit-learn and PyTorch Ecosystem}
\label{sec:tech-sklearn-pytorch}
% ------------------------------------------------------------

The experimental pipeline relies on several additional frameworks:

\begin{itemize}
    \item \textbf{Scikit-learn}: used for data
          splitting (\texttt{train\_test\_split} with stratification), feature
          scaling (\texttt{StandardScaler}), label encoding
          (\texttt{LabelEncoder}), and evaluation metrics (accuracy, F1 score,
          Matthews Correlation Coefficient, ROC-AUC, classification report).
    \item \textbf{PyTorch}: the underlying deep
          learning framework for both TabNet (via \texttt{pytorch-tabnet}) and
          TFT (via \texttt{pytorch-forecasting} and
          \texttt{pytorch-lightning}).
    \item \textbf{PyTorch Lightning}: a high-level
          training framework that abstracts the training loop, gradient
          accumulation, checkpointing, and multi-GPU/distributed training,
          used to train the TFT model.
    \item \textbf{PyTorch Forecasting}:
          a library built on top of PyTorch Lightning that provides pre-built
          time-series architectures including TFT, along with the
          \texttt{TimeSeriesDataSet} abstraction for efficient batched
          sequence generation.
\end{itemize}

\end{document}